% -------------------------------------------------------------
%  RESUMOS
\setlength{\absparsep}{18pt} % ajusta o espaçamento dos parágrafos do resumo
% -------------------------------------------------------------
% ATENÇÃO: o ambiente 'otherlanguage*' deve ser usado para o resumo que não está na
% língua vernácula do trabalho, com a respectiva opção linguística do pacote 'babel'.
% -------------------------------------------------------------
% resumo em PORTUGUÊS (OBRIGATÓRIO)
\begin{resumo}[Resumo]
 \begin{otherlanguage*}{brazil}

   Em um esforço para combater as alarmantes taxas de incidência e mortalidade por câncer de colo do útero no estado do Amazonas, um cenário de saúde pública agravado por barreiras geográficas que dificultam o diagnóstico precoce, este projeto de pesquisa propõe uma investigação aprofundada das bases moleculares da agressividade do Papilomavírus Humano tipo 16 (HPV-16), o genótipo de maior risco oncogênico e prevalência na região. A pesquisa foca-se em como as mutações específicas do HPV-16, circulantes na população amazonense e no mundo, alteram a estrutura e a função de suas oncoproteínas cruciais, a E6 e a E7, potencializando a sua capacidade de induzir a transformação maligna das células hospedeiras. A metodologia empregará uma abordagem computacional, iniciando com a modelagem tridimensional das oncoproteínas E6 e E7, tanto em sua forma selvagem quanto incorporando as variantes genéticas encontradas no Amazonas e no mundo, utilizando para isso o poder da inteligência artificial através da ferramenta AlphaFold. Em seguida, serão realizadas simulações de docking molecular para predizer a interação dessas oncoproteínas virais com seus alvos celulares primários, a proteína supressora de tumor p53, que é inativada pela E6, além de investigar interações com a proteína p16INK4a, um biomarcador fundamental na gravidade das lesões cervicais ocasionadas pelo câncer de colo de útero. Também será realizado uma dinâmica molecular ao longo de 200 ns, que permitirá observar o comportamento dos complexos ao longo do tempo. Além disso, serão  propostos anticorpos neutralizantes do vírus e regiões de epítopos antigênicos, bem como sequência de primers para o diagnóstico. Ao fim,  as sequências de proteínas, epítopos e primers encontradas serão disponibilizadas publicamente ao término deste trabalho no Github, o que possibilitará a síntese, purificação e a implementação dos testes diagnósticos num próximo trabalho.

    \textbf{Palavras-chave}: HPV-16, Amazonas, bioinformática, IA, docking molecular, dinâmica molecular

 \end{otherlanguage*}
\end{resumo}
% -------------------------------------------------------------
% -------------------------------------------------------------
% resumo em INGLÊS (OBRIGATÓRIO)
\begin{resumo}[Abstract]
 \begin{otherlanguage*}{english}
    This is the english abstract.
    
    \textbf{Keywords}: latex. abntex. text editoration.
 \end{otherlanguage*}
\end{resumo}
% -------------------------------------------------------------